\begin{PSbox}[unbreakable]{obj}{اهداف یادگیری}

پس از تکمیل این ماژول، دانشجویان قادر خواهند بود:

\begin{enumerate}
\def\labelenumi{\arabic{enumi}.}
\tightlist
\item
  آزمایش‌های تصادفی را تعریف کنند و فضای نمونه را برای سناریوهای مهندسی شناسایی کنند
\item
  عملیات مجموعه‌ای را برای تعریف رویدادهای پیچیده به کار ببرند
\item
  از اصول موضوعه احتمال برای استخراج قواعد احتمال استفاده کنند
\item
  احتمال‌های شرطی را محاسبه کنند
\item
  استقلال آماری را تعیین کنند
\item
  قضیه بیز را در مسایل استنباط مهندسی به کار ببرند
\item
  از قضیه احتمال کل برای تحلیل چندمسیری استفاده کنند
\item
  آزمایش‌های تصادفی را در \lr{\mbox{MATLAB}} شبیه‌سازی کنند
\end{enumerate}

\end{PSbox}

\PSsection{1.1}{مقدمه: چرا احتمال در مهندسی؟}

\PSsubsection{واقعیت مهندسی}

هیچ سیستم مهندسی در یک دنیای کاملاً قطعی عمل نمی‌کند. موارد زیر را در نظر بگیرید:

\begin{itemize}
\tightlist
\item
  یک سیستم ارتباط دیجیتال بیت‌ها را ارسال می‌کند، اما \textbf{نویز} برخی بیت‌ها را به‌صورت تصادفی خراب می‌کند
\item
  یک خط تولید مقاومت‌ها را تولید می‌کند، اما مقادیر آن‌ها حول مقدار نامی \textbf{تغییر می‌کنند}
\item
  یک حسگر دما را اندازه‌گیری می‌کند، اما هر قرایت دارای \textbf{خطای تصادفی} است
\item
  یک سیگنال بی‌سیم به گیرنده می‌رسد، اما \textbf{محو‌شدگی} باعث تضعیف تصادفی می‌شود
\end{itemize}

\textbf{نظریه احتمال} چارچوب ریاضی لازم برای مدل‌سازی، تحلیل و طراحی سیستم‌هایی را فراهم می‌کند که تحت عدم‌قطعیت کار می‌کنند.

\PSsubsection{پرسش مهندسی}

\begin{PSquote}

«اگر 1000 بیت را روی یک کانال نویزی ارسال کنم، انتظار دارم چند خطا رخ دهد؟ آیا می‌توانم کمتر از 5 خطا را تضمین کنم؟»

\end{PSquote}

برای پاسخ به این پرسش، به موارد زیر نیاز داریم:

\begin{enumerate}
\def\labelenumi{\arabic{enumi}.}
\tightlist
\item
  مدلی از پدیده تصادفی (خطاهای بیت)
\item
  ابزارهای ریاضی برای محاسبه احتمال‌ها
\item
  روش‌هایی برای اتخاذ تصمیم‌های مهندسی بر اساس این احتمال‌ها
\end{enumerate}

\PSsection{1.2}{آزمایش‌های تصادفی}

\begin{PSbox}[unbreakable]{def}{تعریف}

یک \textbf{آزمایش تصادفی} روشی است که:

\begin{enumerate}
\def\labelenumi{\arabic{enumi}.}
\tightlist
\item
  بتوان آن را تحت شرایط یکسان تکرار کرد
\item
  مجموعه‌ای کاملاً تعریف‌شده از خروجی‌های ممکن داشته باشد
\item
  خروجی آن پیش از انجام آزمایش با قطعیت قابل پیش‌بینی نباشد
\end{enumerate}

\end{PSbox}

\PSsubsection{مثال‌های مهندسی}

{\def\LTcaptype{none} % do not increment counter
\begin{longtable}[c]{@{}ccc@{}}
\toprule\noalign{}
\rowcolor{PSnavy}\PSth{آزمایش} & \PSth{زمینه} & \PSth{منبع عدم‌قطعیت} \\
\midrule\noalign{}
\endhead
\bottomrule\noalign{}
\endlastfoot
ارسال یک بیت روی یک کانال نویزی & ارتباطات & نویز حرارتی \\
اندازه‌گیری مقدار یک مقاومت & الکترونیک & تغییرات ساخت \\
شمارش بسته‌های واردشده در 1 ثانیه & شبکه & تصادفی‌بودن ترافیک \\
اندازه‌گیری توان سیگنال دریافتی & ارتباطات بی‌سیم & محوشدگی، تداخل \\
آزمون یک قطعه تا خرابی & قابلیت اطمینان & نقص مواد، فرسایش \\
نمونه‌برداری از ولتاژ خروجی \lr{\mbox{ADC}} & پردازش سیگنال & کوانتیزه‌سازی + نویز \\
\end{longtable}
}

\PSsubsection{تمایز کلیدی}

{\def\LTcaptype{none} % do not increment counter
\begin{longtable}[c]{@{}
  >{\centering\arraybackslash}p{(0.965\linewidth - 2\tabcolsep) * \real{0.6364}}
  >{\centering\arraybackslash}p{(0.965\linewidth - 2\tabcolsep) * \real{0.3636}}@{}}
\toprule\noalign{}
\rowcolor{PSnavy}\PSth{قطعی} & \PSth{تصادفی} \\
\midrule\noalign{}
\endhead
\bottomrule\noalign{}
\endlastfoot
خروجی یک گیت منطقی با ورودی‌های معلوم & خروجی یک گیت منطقی با ورودی‌های نویزی \\
مقاومت محاسبه‌شده از روی کد رنگی & مقاومت واقعی اندازه‌گیری‌شده \\
تلفات مسیر در فضای آزاد در فاصله معلوم & توان دریافتی واقعی همراه با محوشدگی \\
\end{longtable}
}

\PSsection{1.3}{فضاهای نمونه}

\begin{PSbox}[unbreakable]{def}{تعریف}

\textbf{فضای نمونه} \lr{$S$} (یا \lr{$\Omega$}) مجموعه تمام خروجی‌های ممکن یک آزمایش تصادفی است.

\end{PSbox}

\PSsubsection{انواع فضاهای نمونه}

\textbf{گسسته (شمارا):}

\begin{itemize}
\tightlist
\item
  ارسال یک بیت: \lr{$S = \{0,\, 1\}$} (دریافت صحیح یا دریافت همراه با خطا)
\item
  تعداد خطاها در یک بسته \lr{$n$} بیتی: \lr{$S = \{0,\, 1,\, 2,\, \ldots,\, n\}$}
\item
  تعداد ورودها به یک سرور: \lr{$S = \{0,\, 1,\, 2,\, 3,\, \ldots \}$}
\end{itemize}

\textbf{پیوسته (ناشمارا):}

\begin{itemize}
\tightlist
\item
  ولتاژ اندازه‌گیری‌شده: \lr{$S = \mathbb{R}$} (یا یک بازه عملی مانند \lr{$[0,\, 5] V$})
\item
  زمان تا خرابی قطعه: \lr{$S = [0,\, \infty)$}
\item
  فاز سیگنال دریافتی: \lr{$S = [0,\, 2\pi)$}
\end{itemize}

\begin{PSbox}[unbreakable]{ex}{مثال مهندسی: سیستم ارتباطی}

یک فرستنده یک بیت (0 یا 1) را روی یک کانال نویزی ارسال می‌کند. گیرنده تصمیم‌گیری می‌کند.

\PSfigure{m01-sample-space}

\end{PSbox}

\PSsection{1.4}{رویدادها}

\begin{PSbox}[unbreakable]{def}{تعریف}

یک \textbf{رویداد} زیرمجموعه‌ای از فضای نمونه است. رویداد \textbf{رخ می‌دهد} اگر خروجی آزمایش به آن زیرمجموعه تعلق داشته باشد.

\end{PSbox}

\PSsubsection{انواع رویدادها}

\begin{itemize}
\tightlist
\item
  \textbf{رویداد ساده}: دقیقاً شامل یک خروجی است. \lr{$E = \{(0,\,1)\}$}
\item
  \textbf{رویداد مرکب}: شامل چندین خروجی است. \lr{$E = \{(0,\,1),\, (1,\,0)\}$} (هر خطایی)
\item
  \textbf{رویداد قطعی}: کل فضای نمونه \lr{$S$} (همیشه رخ می‌دهد)
\item
  \textbf{رویداد ناممکن}: مجموعه تهی \lr{$\emptyset$} (هرگز رخ نمی‌دهد)
\end{itemize}

\PSsubsection{عملیات مجموعه‌ای روی رویدادها}

{\def\LTcaptype{none} % do not increment counter
\begin{longtable}[c]{@{}
  >{\centering\arraybackslash}p{(0.965\linewidth - 6\tabcolsep) * \real{0.1930}}
  >{\centering\arraybackslash}p{(0.965\linewidth - 6\tabcolsep) * \real{0.1754}}
  >{\centering\arraybackslash}p{(0.965\linewidth - 6\tabcolsep) * \real{0.1579}}
  >{\centering\arraybackslash}p{(0.965\linewidth - 6\tabcolsep) * \real{0.4737}}@{}}
\toprule\noalign{}
\rowcolor{PSnavy}\PSth{عملیات} & \PSth{نمادگذاری} & \PSth{معنا} & \PSth{تفسیر مهندسی} \\
\midrule\noalign{}
\endhead
\bottomrule\noalign{}
\endlastfoot
اجتماع & \lr{$A \cup B$} & \lr{$A$} یا \lr{$B$} (یا هر دو) رخ می‌دهد & هر یک از حالت‌های خرابی \lr{$A$} یا \lr{$B$} موجب فعال‌شدن هشدار می‌شود \\
اشتراک & \lr{$A \cap B$} & \lr{$A$} و \lr{$B$} هر دو رخ می‌دهند & هر دو حسگر سیگنال را تشخیص می‌دهند \\
متمم & \lr{\mbox{$A^{c}$ or Ā}} & \lr{$A$} رخ نمی‌دهد & قطعه خراب نمی‌شود \\
تفاضل & \lr{$A - B = A \cap B^{c}$} & \lr{$A$} رخ می‌دهد ولی \lr{$B$} رخ نمی‌دهد & حسگر \lr{$A$} فعال می‌شود ولی حسگر \lr{$B$} فعال نمی‌شود \\
ناسازگار متقابل & \lr{$A \cap B = \emptyset$} & \lr{$A$} و \lr{$B$} نمی‌توانند هر دو رخ دهند & قطعه نمی‌تواند هم سالم و هم خراب باشد \\
\end{longtable}
}

\begin{PSbox}[unbreakable]{ex}{مثال مهندسی: عیوب ساخت}

بازرسی یک برد مدار سه نوع عیب را بررسی می‌کند:

\begin{itemize}
\tightlist
\item
  \lr{$A$} = عیب پل لحیم
\item
  \lr{$B$} = قطعه مفقود
\item
  \lr{$C$} = مقدار نادرست قطعه
\end{itemize}

رویدادها:

\begin{itemize}
\tightlist
\item
  \lr{$A \cup B \cup C$} = «برد دست‌کم یک عیب دارد»
\item
  \lr{$A \cap B$} = «برد هم عیب پل لحیم و هم یک قطعه مفقود دارد»
\item
  \lr{$(A \cup B \cup C)^{c}$} = «برد بدون عیب است»
\end{itemize}

\end{PSbox}

\PSsection{1.5}{اصول موضوعه احتمال}

\PSsubsection{اصول موضوعه کلموگروف}

برای یک فضای نمونه \lr{$S$} که رویدادهایی روی آن تعریف شده‌اند، تابع احتمال \lr{$P$} شرایط زیر را برآورده می‌کند:

\textbf{اصل موضوعه 1 (نامنفی‌بودن):}\PSdisplay{P(A) \geq 0 \quad \PSmtext{for any event } A}

\textbf{اصل موضوعه 2 (نرمال‌سازی):}\PSdisplay{P(S) = 1}

\textbf{اصل موضوعه 3 (جمع‌پذیری شمارا):}\PSbr{}برای رویدادهای ناسازگار متقابل \lr{$A_{1}$}، \lr{$A_{2}$}، \lr{$A_{3}$}، \lr{…}:\PSdisplay{P(A_1 \cup A_2 \cup A_3 \cup \cdots) = P(A_1) + P(A_2) + P(A_3) + \cdots}

\PSsubsection{ویژگی‌های استخراج‌شده از اصول موضوعه}

\textbf{ویژگی 1: قاعده متمم}\PSdisplay{P(A^c) = 1 - P(A)}

\emph{اثبات:} \lr{$S = A \cup A^{c}$} و \lr{$A \cap A^{c} = \emptyset$}، بنابراین \lr{$P(S) = P(A) + P(A^{c}) = 1$}.

\emph{معنای مهندسی:} \lr{$P(\text{\rl{عملکرد قطعه}}) = 1 - P(\text{\rl{خرابی قطعه}})$}

\textbf{ویژگی 2: رویداد ناممکن}\PSdisplay{P(\emptyset) = 0}

\textbf{ویژگی 3: یکنوایی}\PSbr{}اگر \lr{$A \subseteq B$}، آنگاه \lr{$P(A) \le P(B)$}

\textbf{ویژگی 4: کران}\PSdisplay{0 \leq P(A) \leq 1}

\textbf{ویژگی 5: شمول و عدم‌شمول (دو رویداد)}\PSdisplay{P(A \cup B) = P(A) + P(B) - P(A \cap B)}

\emph{معنای مهندسی:} هنگام محاسبه احتمال «خرابی \lr{$A$} یا خرابی \lr{\mbox{B}}»، باید هم‌پوشانی‌ای را که در آن هر دو به‌طور هم‌زمان خراب می‌شوند، کم کنیم.

\PSsection{1.6}{قواعد احتمال}

\PSsubsection{قاعده جمع}

برای هر دو رویداد:\PSdisplay{P(A \cup B) = P(A) + P(B) - P(A \cap B)}

برای رویدادهای ناسازگار متقابل \lr{$(A \cap B = \emptyset)$}:\PSdisplay{P(A \cup B) = P(A) + P(B)}

\PSsubsection{قاعده متمم}

\PSdisplay{P(A) = 1 - P(A^c)}

این قاعده اغلب ساده‌ترین روش برای محاسبه احتمال‌هاست:

\begin{PSquote}

\lr{$P(\text{\rl{دست‌کم یک خطا در 100 بیت}}) = 1 - P(\text{\rl{صفر خطا در 100 بیت}})$}

\end{PSquote}

\PSsubsection{شمول و عدم‌شمول (سه رویداد)}

\PSdisplay{P(A \cup B \cup C) = P(A) + P(B) + P(C) - P(A \cap B) - P(A \cap C) - P(B \cap C) + P(A \cap B \cap C)}

\begin{PSbox}[unbreakable]{ex}{مثال مهندسی: افزونگی شبکه}

یک مرکز داده دارای سه پیوند شبکه مستقل است. هر پیوند با احتمال \PSnum{0.05} خراب می‌شود.

\begin{itemize}
\tightlist
\item
  \lr{$A$} = خرابی پیوند 1، \lr{$B$} = خرابی پیوند 2، \lr{$C$} = خرابی پیوند 3
\item
  \lr{$P(\text{\rl{دست‌کم یکی خراب می‌شود}}) = P(A \cup B \cup C)$}
\item
  با استفاده از شمول و عدم‌شمول:
\end{itemize}

\lr{$P(A \cup B \cup C) = 3(0.05) - 3(0.05)^{2} + (0.05)^{3} = 0.15 - 0.0075 + 0.000125 = 0.142625$}

به‌صورت جایگزین: \lr{$P(\text{\rl{دست‌کم یکی خراب می‌شود}}) = 1 - P(\text{\rl{هیچ‌کدام خراب نمی‌شوند}}) = 1 - (0.95)^{3} = 0.142625$ \PSyes{}}

\end{PSbox}

\PSsection{1.7}{احتمال شرطی}

\begin{PSbox}[unbreakable]{def}{تعریف}

\textbf{احتمال شرطی} رویداد \lr{$A$} با توجه به اینکه رویداد \lr{$B$} رخ داده است:

\PSdisplay{P(A|B) = \frac{P(A \cap B)}{P(B)}, \quad P(B) > 0}

\end{PSbox}

\PSsubsection{تفسیر}

احتمال شرطی دانش ما را \textbf{به‌روزرسانی} می‌کند. هنگامی که می‌دانیم \lr{$B$} رخ داده است، فضای نمونه مؤثر از \lr{$S$} به \lr{$B$} کاهش می‌یابد و می‌پرسیم: «درون \lr{$B$}، احتمال \lr{$A$} چقدر است؟»

\PSsubsection{تمایز کلیدی}

{\def\LTcaptype{none} % do not increment counter
\begin{longtable}[c]{@{}
  >{\centering\arraybackslash}p{(0.965\linewidth - 4\tabcolsep) * \real{0.3214}}
  >{\centering\arraybackslash}p{(0.965\linewidth - 4\tabcolsep) * \real{0.3571}}
  >{\centering\arraybackslash}p{(0.965\linewidth - 4\tabcolsep) * \real{0.3214}}@{}}
\toprule\noalign{}
\rowcolor{PSnavy}\PSth{مفهوم} & \PSth{نمادگذاری} & \PSth{معنا} \\
\midrule\noalign{}
\endhead
\bottomrule\noalign{}
\endlastfoot
احتمال \lr{$A$} & \lr{$P(A)$} & احتمال \lr{$A$} بدون اطلاعات اضافی \\
احتمال شرطی & \lr{$P(A \mid B)$} & احتمال \lr{$A$} \textbf{با توجه به اینکه} \lr{$B$} رخ داده است \\
احتمال مشترک & \lr{$P(A \cap B)$} & احتمال اینکه \textbf{هر دو} \lr{$A$} و \lr{$B$} رخ دهند \\
\end{longtable}
}

این کمیت‌ها با یکدیگر مرتبط‌اند: \lr{$P(A \cap B) = P(A \mid B) \cdot P(B) = P(B \mid A) \cdot P(A)$}

\begin{PSbox}[unbreakable]{ex}{مثال مهندسی: نظریه آشکارسازی}

یک سامانه راداری هواپیماها را آشکار می‌کند:

\begin{itemize}
\tightlist
\item
  \lr{$H_{1}$} = هواپیما وجود دارد (هدف)
\item
  \lr{$H_{0}$} = هواپیما وجود ندارد (بدون هدف)
\item
  \lr{$D$} = آشکارساز می‌گوید «هدف وجود دارد»
\end{itemize}

احتمال‌های شرطی مهم:

\begin{itemize}
\tightlist
\item
  \lr{$P(D \mid H_{1})$} = \textbf{احتمال آشکارسازی} (حساسیت)
\item
  \lr{$P(D \mid H_{0})$} = \textbf{احتمال هشدار کاذب}
\item
  \lr{$P(D^{c} \mid H_{1})$} = \textbf{احتمال عدم‌آشکارسازی} \lr{$= 1 - P(D \mid H_{1})$}
\end{itemize}

فرض کنید: \lr{$P(H_{1}) = 0.01$}، \lr{$P(D \mid H_{1}) = 0.95$}، \lr{$P(D \mid H_{0}) = 0.02$}

\lr{$P(\text{\rl{هدف واقعاً وجود دارد}}\ \mid \text{\rl{آشکارساز وجود هدف را اعلام می‌کند}})$} چقدر است؟

این مسیله به \textbf{قضیه بیز} نیاز دارد (بخش \PSnum{1.9).}

\end{PSbox}

\begin{PSbox}[unbreakable]{ex}{مثال مهندسی: کانال ارتباطی}

یک کانال متقارن دودویی دارای احتمال وارونگی \lr{$p = 0.01$} است:

\begin{itemize}
\tightlist
\item
  \lr{$P(\text{\rl{دریافت 1}}\ \mid \text{\rl{ارسال 0}}) = p = 0.01$}
\item
  \lr{$P(\text{\rl{دریافت 0}}\ \mid \text{\rl{ارسال 1}}) = p = 0.01$}
\item
  \lr{$P(\text{\rl{دریافت 0}}\ \mid \text{\rl{ارسال 0}}) = 1 - p = 0.99$}
\item
  \lr{$P(\text{\rl{دریافت 1}}\ \mid \text{\rl{ارسال 1}}) = 1 - p = 0.99$}
\end{itemize}

اگر بیت‌ها با احتمال مساوی 0 یا 1 باشند:

\begin{itemize}
\tightlist
\item
  \lr{$P(\text{\rl{ارسال 0}}\ \cap \text{\rl{دریافت 1}}) = P(\text{\rl{دریافت 1}}\ \mid \text{\rl{ارسال 0}}) \cdot P(\text{\rl{ارسال 0}}) = 0.01 \times 0.5 = 0.005$}
\item
  \lr{$P(\text{\rl{خطا}}) = P(\text{\rl{ارسال 0}}\ \cap \text{\rl{دریافت 1}}) + P(\text{\rl{ارسال 1}}\ \cap \text{\rl{دریافت 0}}) = 0.005 + 0.005 = 0.01$}
\end{itemize}

\end{PSbox}

\PSsubsection{ویژگی‌های احتمال شرطی}

احتمال شرطی خود یک معیار احتمال معتبر است:

\begin{enumerate}
\def\labelenumi{\arabic{enumi}.}
\tightlist
\item
  \lr{$P(A \mid B) \ge 0$}
\item
  \lr{$P(S \mid B) = 1$}
\item
  \lr{$P(A_{1} \cup A_{2} \mid B) = P(A_{1} \mid B) + P(A_{2} \mid B)$} اگر \lr{$A_{1} \cap A_{2} = \emptyset$}
\end{enumerate}

\PSsubsection{قاعده ضرب (قاعده زنجیره‌ای)}

\PSdisplay{P(A \cap B) = P(A|B) \cdot P(B) = P(B|A) \cdot P(A)}

برای سه رویداد:\PSdisplay{P(A \cap B \cap C) = P(A) \cdot P(B|A) \cdot P(C|A \cap B)}

\PSsection{1.8}{استقلال}

\begin{PSbox}[unbreakable]{def}{تعریف}

رویدادهای \lr{$A$} و \lr{$B$} از نظر \textbf{آماری مستقل} هستند اگر:

\PSdisplay{P(A \cap B) = P(A) \cdot P(B)}

به‌طور معادل: \lr{$P(A \mid B) = P(A)$} \lr{—} دانستن اینکه \lr{$B$} رخ داده است، هیچ اطلاعاتی درباره \lr{$A$} نمی‌دهد.

\end{PSbox}

\PSsubsection{اهمیت مهندسی}

استقلال به این معناست که \textbf{یک رویداد هیچ اطلاعاتی درباره رویداد دیگر فراهم نمی‌کند}. این یک فرض مدل‌سازی است که اغلب به کار می‌بریم:

\begin{itemize}
\tightlist
\item
  نمونه‌های نویز در زمان‌های مختلف مستقل هستند
\item
  خطاهای بیت در شکاف‌های زمانی مختلف مستقل هستند (کانال بدون حافظه)
\item
  خرابی قطعات در زیرسامانه‌های مختلف مستقل است
\end{itemize}

\PSsubsection{استقلال زوجی در برابر استقلال متقابل}

برای سه رویداد \lr{$A$}، \lr{$B$}، \lr{$C$}:

\textbf{استقلال زوجی} (هر سه زوج):

\begin{itemize}
\tightlist
\item
  \lr{$P(A \cap B) = P(A)P(B)$}
\item
  \lr{$P(A \cap C) = P(A)P(C)$}
\item
  \lr{$P(B \cap C) = P(B)P(C)$}
\end{itemize}

\textbf{استقلال متقابل} (تمام ترکیب‌ها):

\begin{itemize}
\tightlist
\item
  تمام شرایط استقلال زوجی بالا، و
\item
  \lr{$P(A \cap B \cap C) = P(A)P(B)P(C)$}
\end{itemize}

\begin{PSquote}

\lr{\mbox{\PSwarn{}}} استقلال زوجی، استقلال متقابل را نتیجه نمی‌دهد!

\end{PSquote}

\begin{PSbox}[unbreakable]{ex}{مثال مهندسی: خرابی مستقل قطعات}

یک سیستم دارای سه ماژول مستقل است که احتمال خرابی هر یک \PSnum{0.1} است.

\lr{$P(\text{\rl{هر سه خراب شوند}}) = P(A)P(B)P(C) = (0.1)^{3} = 0.001$}

\lr{$P(\text{\rl{سیستم کار می‌کند}}) = P(\text{\rl{دست‌کم یکی کار می‌کند}}) = 1 - P(\text{\rl{همه خراب شوند}}) = 1 - 0.001 = 0.999$}

این محاسبه فقط به این دلیل معتبر است که خرابی‌ها مستقل هستند.

\end{PSbox}

\PSsubsection{آزمون استقلال}

برای بررسی استقلال از روی داده‌ها:

\begin{enumerate}
\def\labelenumi{\arabic{enumi}.}
\tightlist
\item
  \lr{$P(A)$}، \lr{$P(B)$} و \lr{$P(A \cap B)$} را به‌صورت تجربی محاسبه کنید
\item
  بررسی کنید که آیا \lr{$P(A \cap B) \approx P(A) \cdot P(B)$} است
\end{enumerate}

اگر تقریباً برابر نباشند، رویدادها \textbf{وابسته} هستند.

\PSsubsection{برداشت نادرست رایج}

\begin{PSquote}

\lr{\mbox{\PSno{}}} «رویدادهای ناسازگار متقابل مستقل هستند»

\end{PSquote}

این گزاره \textbf{نادرست} است. اگر \lr{$A \cap B = \emptyset$} و \lr{$P(A) > 0$}، \lr{$P(B) > 0$}، آنگاه:

\begin{itemize}
\tightlist
\item
  \lr{$P(A \cap B) = 0$}
\item
  \lr{$P(A)P(B) > 0$}
\item
  بنابراین \lr{$P(A \cap B) \ne P(A)P(B)$}
\end{itemize}

رویدادهای ناسازگار متقابل در واقع \textbf{به‌شدت وابسته} هستند \lr{—} اگر یکی رخ دهد، دیگری نمی‌تواند رخ دهد.

\PSsection{1.9}{قضیه بیز}

\PSsubsection{استخراج}

از تعریف احتمال شرطی:

\PSdisplay{P(A|B) = \frac{P(A \cap B)}{P(B)} = \frac{P(B|A) \cdot P(A)}{P(B)}}

\PSsubsection{بیان کامل}

اگر \lr{$B_{1}$}، \lr{$B_{2}$}، \lr{…}، \lr{$B_{n}$} فضای نمونه را افراز کنند (ناسازگار متقابل و جامع باشند):

\PSdisplay{P(B_i|A) = \frac{P(A|B_i) \cdot P(B_i)}{\sum_{j=1}^{n} P(A|B_j) \cdot P(B_j)}}

\PSsubsection{اصطلاحات}

{\def\LTcaptype{none} % do not increment counter
\begin{longtable}[c]{@{}
  >{\centering\arraybackslash}p{(0.965\linewidth - 4\tabcolsep) * \real{0.1818}}
  >{\centering\arraybackslash}p{(0.965\linewidth - 4\tabcolsep) * \real{0.1818}}
  >{\centering\arraybackslash}p{(0.965\linewidth - 4\tabcolsep) * \real{0.6364}}@{}}
\toprule\noalign{}
\rowcolor{PSnavy}\PSth{اصطلاح} & \PSth{نام} & \PSth{معنای مهندسی} \\
\midrule\noalign{}
\endhead
\bottomrule\noalign{}
\endlastfoot
\lr{$P(B_{i})$} & احتمال پیشین & آنچه پیش از مشاهده داده‌ها باور داشتیم \\
\lr{$P(A \mid B_{i})$} & درست‌نمایی & با توجه به هر فرضیه، مشاهده با چه احتمالی رخ می‌دهد \\
\lr{$P(B_{i} \mid A)$} & احتمال پسین & باور به‌روزشده پس از مشاهده داده‌ها \\
\lr{$P(A)$} & شواهد & احتمال کل مشاهده \\
\end{longtable}
}

\begin{PSbox}[unbreakable]{ex}{مثال مهندسی: تشخیص خطا}

یک برد مدار می‌تواند در یکی از سه زیرسامانه دارای خطا باشد:

\begin{itemize}
\tightlist
\item
  \lr{$B_{1}$} = خطای منبع تغذیه، \lr{$P(B_{1}) = 0.3$}
\item
  \lr{$B_{2}$} = خطای پردازنده، \lr{$P(B_{2}) = 0.5$}
\item
  \lr{$B_{3}$} = خطای حافظه، \lr{$P(B_{3}) = 0.2$}
\end{itemize}

یک آزمون تشخیصی علامت \lr{$A$} (گرم‌شدن بیش از حد) را ایجاد می‌کند:

\begin{itemize}
\tightlist
\item
  \lr{$P(A \mid B_{1}) = 0.8$} (خطاهای منبع تغذیه معمولاً باعث گرم‌شدن بیش از حد می‌شوند)
\item
  \lr{$P(A \mid B_{2}) = 0.3$} (خطاهای پردازنده گاهی باعث گرم‌شدن بیش از حد می‌شوند)
\item
  \lr{$P(A \mid B_{3}) = 0.1$} (خطاهای حافظه به‌ندرت باعث گرم‌شدن بیش از حد می‌شوند)
\end{itemize}

\textbf{پرسش:} با توجه به مشاهده گرم‌شدن بیش از حد، محتمل‌ترین خطا کدام است؟

\textbf{راه‌حل:}

\lr{$P(A) = P(A \mid B_{1})P(B_{1}) + P(A \mid B_{2})P(B_{2}) + P(A \mid B_{3})P(B_{3})$}\PSbr{}\lr{$= (0.8)(0.3) + (0.3)(0.5) + (0.1)(0.2)$}\PSbr{}\lr{$= 0.24 + 0.15 + 0.02 = 0.41$}

\lr{$P(B_{1} \mid A) = \frac{(0.8)(0.3)}{0.41} = 0.585$} (محتمل‌ترین!)\PSbr{}\lr{$\displaystyle P(B_{2} \mid A) = \frac{(0.3)(0.5)}{0.41} = 0.366$}\PSbr{}\lr{$\displaystyle P(B_{3} \mid A) = \frac{(0.1)(0.2)}{0.41} = 0.049$}

\textbf{تصمیم مهندسی:} ابتدا منبع تغذیه را بررسی کنید.

\end{PSbox}

\begin{PSbox}[unbreakable]{ex}{مثال مهندسی: آشکارسازی با بیز}

به مثال رادار از بخش \PSnum{1.7} بازمی‌گردیم:

\begin{itemize}
\tightlist
\item
  \lr{$P(H_{1}) = 0.01,\, P(D \mid H_{1}) = 0.95,\, P(D \mid H_{0}) = 0.02$}
\end{itemize}

\lr{$P(D) = P(D \mid H_{1})P(H_{1}) + P(D \mid H_{0})P(H_{0}) = (0.95)(0.01) + (0.02)(0.99) = 0.0095 + 0.0198 = 0.0293$}

\lr{$\displaystyle P(H_{1} \mid D) = \frac{(0.95)(0.01)}{0.0293} = 0.324$}

\textbf{تفسیر:} حتی با یک آشکارساز خوب (95\% آشکارسازی، 2\% هشدار کاذب)، هنگامی که اهداف نادر هستند (1\%)، بیشتر آشکارسازی‌ها هشدار کاذب‌اند! فقط \PSnum{32.4\%} از هشدارها با اهداف واقعی مطابقت دارند.

\textbf{درس مهندسی:} طراحی سیستم باید \textbf{احتمال‌های پیشین} را در نظر بگیرد، نه فقط عملکرد آشکارساز را.

\end{PSbox}

\PSsection{1.10}{قضیه احتمال کل}

\PSsubsection{بیان}

اگر \lr{$B_{1}$}، \lr{$B_{2}$}، \lr{…}، \lr{$B_{n}$} یک \textbf{افراز} از \lr{$S$} تشکیل دهند (ناسازگار متقابل و جامع):

\PSdisplay{P(A) = \sum_{i=1}^{n} P(A|B_i) \cdot P(B_i)}

\PSsubsection{تفسیر}

برای یافتن احتمال کل رویداد \lr{$A$}، آن را به تمام «مسیرهای» ممکن منتهی به \lr{$A$} از طریق افراز تجزیه کنید.

\begin{PSbox}[unbreakable]{ex}{مثال مهندسی: ارتباط چندمسیری}

یک بسته می‌تواند از طریق سه مسیر مختلف مسیریابی شود:

\begin{itemize}
\tightlist
\item
  مسیر 1: در 50\% مواقع انتخاب می‌شود، احتمال خطا \PSnum{0.01}
\item
  مسیر 2: در 30\% مواقع انتخاب می‌شود، احتمال خطا \PSnum{0.05}
\item
  مسیر 3: در 20\% مواقع انتخاب می‌شود، احتمال خطا \PSnum{0.10}
\end{itemize}

\lr{$P(\text{\rl{خطای بسته}}) = P(\PSmtext{error}\ \mid \PSmtext{path 1})P(\PSmtext{path 1}) + P(\PSmtext{error}\ \mid \PSmtext{path 2})P(\PSmtext{path 2}) + P(\PSmtext{error}\ \mid \PSmtext{path 3})P(\PSmtext{path 3})$}\PSbr{}\lr{$= (0.01)(0.50) + (0.05)(0.30) + (0.10)(0.20)$}\PSbr{}\lr{$= 0.005 + 0.015 + 0.020 = 0.040$}

\textbf{نرخ خطای کلی: 4\%}

\end{PSbox}

\begin{PSbox}[unbreakable]{ex}{مثال مهندسی: تولید با تأمین‌کنندگان متعدد}

یک شرکت خازن‌ها را از دو تأمین‌کننده تهیه می‌کند:

\begin{itemize}
\tightlist
\item
  تأمین‌کننده \lr{\mbox{$A$: 60\%}} خازن‌ها را تأمین می‌کند، نرخ عیب 2\%
\item
  تأمین‌کننده \lr{\mbox{$B$: 40\%}} خازن‌ها را تأمین می‌کند، نرخ عیب 5\%
\end{itemize}

\lr{$P(\text{\rl{معیوب}}) = P(\text{\rl{معیوب}}\ \mid A)P(A) + P(\text{\rl{معیوب}}\ \mid B)P(B)$}\PSbr{}\lr{$= (0.02)(0.60) + (0.05)(0.40)$}\PSbr{}\lr{$= 0.012 + 0.020 = 0.032$}

اگر یک خازن معیوب پیدا شود، کدام تأمین‌کننده احتمالاً آن را تولید کرده است؟

\lr{$\displaystyle P(A \mid \text{\rl{معیوب}}) = \frac{(0.02)(0.60)}{0.032} = 0.375$}\PSbr{}\lr{$\displaystyle P(B \mid \text{\rl{معیوب}}) = \frac{(0.05)(0.40)}{0.032} = 0.625$}

با وجود تأمین تعداد کمتری از واحدها، احتمال اینکه تأمین‌کننده \lr{$B$} مسیول یک عیب مشخص باشد بیشتر است.

\end{PSbox}

\PSsection{1.11}{مثال‌های \lr{\mbox{MATLAB}}}

\begin{PSbox}[unbreakable]{ex}{مثال 1: شبیه‌سازی پرتاب سکه (آزمایش تصادفی)}

\tcbset{PScodemode/.style={unbreakable}}

\begin{Shaded}
\begin{Highlighting}[]
\CommentTok{\%\% Simulating a Random Experiment: Biased Coin}
\CommentTok{\% A communication channel has bit error probability p = 0.1}
\CommentTok{\% Simulate transmitting N bits and count errors}

\VariableTok{p} \OperatorTok{=} \FloatTok{0.1}\OperatorTok{;}          \CommentTok{\% Bit error probability}
\VariableTok{N} \OperatorTok{=} \FloatTok{10000}\OperatorTok{;}        \CommentTok{\% Number of bits transmitted}

\CommentTok{\% Generate random outcomes: 1 = error, 0 = correct}
\VariableTok{errors} \OperatorTok{=} \VariableTok{rand}\NormalTok{(}\FloatTok{1}\OperatorTok{,} \VariableTok{N}\NormalTok{) }\OperatorTok{\textless{}} \VariableTok{p}\OperatorTok{;}

\CommentTok{\% Empirical probability of error}
\VariableTok{p\_empirical} \OperatorTok{=} \VariableTok{sum}\NormalTok{(}\VariableTok{errors}\NormalTok{) }\OperatorTok{/} \VariableTok{N}\OperatorTok{;}

\VariableTok{fprintf}\NormalTok{(}\SpecialStringTok{\textquotesingle{}Theoretical P(error) = \%.4f\textbackslash{}n\textquotesingle{}}\OperatorTok{,} \VariableTok{p}\NormalTok{)}\OperatorTok{;}
\VariableTok{fprintf}\NormalTok{(}\SpecialStringTok{\textquotesingle{}Empirical P(error)   = \%.4f (from \%d trials)\textbackslash{}n\textquotesingle{}}\OperatorTok{,} \VariableTok{p\_empirical}\OperatorTok{,} \VariableTok{N}\NormalTok{)}\OperatorTok{;}
\VariableTok{fprintf}\NormalTok{(}\SpecialStringTok{\textquotesingle{}Difference           = \%.4f\textbackslash{}n\textquotesingle{}}\OperatorTok{,} \VariableTok{abs}\NormalTok{(}\VariableTok{p} \OperatorTok{{-}} \VariableTok{p\_empirical}\NormalTok{))}\OperatorTok{;}
\end{Highlighting}
\end{Shaded}

\end{PSbox}

\begin{PSbox}[unbreakable]{ex}{مثال 2: راستی‌آزمایی قاعده جمع}

\tcbset{PScodemode/.style={unbreakable}}

\begin{Shaded}
\begin{Highlighting}[]
\CommentTok{\%\% Verifying Inclusion{-}Exclusion with Simulation}
\CommentTok{\% Two independent failure modes:}
\CommentTok{\% A = overheating (P = 0.05)}
\CommentTok{\% B = power surge (P = 0.03)}

\VariableTok{N} \OperatorTok{=} \FloatTok{100000}\OperatorTok{;}       \CommentTok{\% Number of trials}
\VariableTok{pA} \OperatorTok{=} \FloatTok{0.05}\OperatorTok{;}
\VariableTok{pB} \OperatorTok{=} \FloatTok{0.03}\OperatorTok{;}

\CommentTok{\% Generate independent failures}
\VariableTok{A} \OperatorTok{=} \VariableTok{rand}\NormalTok{(}\FloatTok{1}\OperatorTok{,} \VariableTok{N}\NormalTok{) }\OperatorTok{\textless{}} \VariableTok{pA}\OperatorTok{;}
\VariableTok{B} \OperatorTok{=} \VariableTok{rand}\NormalTok{(}\FloatTok{1}\OperatorTok{,} \VariableTok{N}\NormalTok{) }\OperatorTok{\textless{}} \VariableTok{pB}\OperatorTok{;}

\CommentTok{\% Empirical probabilities}
\VariableTok{pA\_emp} \OperatorTok{=} \VariableTok{mean}\NormalTok{(}\VariableTok{A}\NormalTok{)}\OperatorTok{;}
\VariableTok{pB\_emp} \OperatorTok{=} \VariableTok{mean}\NormalTok{(}\VariableTok{B}\NormalTok{)}\OperatorTok{;}
\VariableTok{pAB\_emp} \OperatorTok{=} \VariableTok{mean}\NormalTok{(}\VariableTok{A} \OperatorTok{\&} \VariableTok{B}\NormalTok{)}\OperatorTok{;}           \CommentTok{\% Both occur}
\VariableTok{pAuB\_emp} \OperatorTok{=} \VariableTok{mean}\NormalTok{(}\VariableTok{A} \OperatorTok{|} \VariableTok{B}\NormalTok{)}\OperatorTok{;}          \CommentTok{\% At least one occurs}

\CommentTok{\% Theoretical (independent)}
\VariableTok{pAB\_theory} \OperatorTok{=} \VariableTok{pA} \OperatorTok{*} \VariableTok{pB}\OperatorTok{;}
\VariableTok{pAuB\_theory} \OperatorTok{=} \VariableTok{pA} \OperatorTok{+} \VariableTok{pB} \OperatorTok{{-}} \VariableTok{pA}\OperatorTok{*}\VariableTok{pB}\OperatorTok{;}

\VariableTok{fprintf}\NormalTok{(}\SpecialStringTok{\textquotesingle{}P(A ∩ B): Theory = \%.5f, Empirical = \%.5f\textbackslash{}n\textquotesingle{}}\OperatorTok{,} \VariableTok{pAB\_theory}\OperatorTok{,} \VariableTok{pAB\_emp}\NormalTok{)}\OperatorTok{;}
\VariableTok{fprintf}\NormalTok{(}\SpecialStringTok{\textquotesingle{}P(A ∪ B): Theory = \%.5f, Empirical = \%.5f\textbackslash{}n\textquotesingle{}}\OperatorTok{,} \VariableTok{pAuB\_theory}\OperatorTok{,} \VariableTok{pAuB\_emp}\NormalTok{)}\OperatorTok{;}
\end{Highlighting}
\end{Shaded}

\end{PSbox}

\begin{PSbox}[unbreakable]{ex}{مثال 3: شبیه‌سازی احتمال شرطی}

\tcbset{PScodemode/.style={unbreakable}}

\begin{Shaded}
\begin{Highlighting}[]
\CommentTok{\%\% Conditional Probability: Binary Symmetric Channel}
\CommentTok{\% P(error | sent 0) = P(error | sent 1) = p = 0.1}

\VariableTok{p\_crossover} \OperatorTok{=} \FloatTok{0.1}\OperatorTok{;}
\VariableTok{N} \OperatorTok{=} \FloatTok{100000}\OperatorTok{;}

\CommentTok{\% Randomly generate transmitted bits (equally likely 0 or 1)}
\VariableTok{transmitted} \OperatorTok{=} \VariableTok{randi}\NormalTok{([}\FloatTok{0}\OperatorTok{,} \FloatTok{1}\NormalTok{]}\OperatorTok{,} \FloatTok{1}\OperatorTok{,} \VariableTok{N}\NormalTok{)}\OperatorTok{;}

\CommentTok{\% Channel: flip each bit with probability p\_crossover}
\VariableTok{noise} \OperatorTok{=} \VariableTok{rand}\NormalTok{(}\FloatTok{1}\OperatorTok{,} \VariableTok{N}\NormalTok{) }\OperatorTok{\textless{}} \VariableTok{p\_crossover}\OperatorTok{;}
\VariableTok{received} \OperatorTok{=} \VariableTok{xor}\NormalTok{(}\VariableTok{transmitted}\OperatorTok{,} \VariableTok{noise}\NormalTok{)}\OperatorTok{;}

\CommentTok{\% Compute conditional probabilities}
\VariableTok{sent\_0\_idx} \OperatorTok{=}\NormalTok{ (}\VariableTok{transmitted} \OperatorTok{==} \FloatTok{0}\NormalTok{)}\OperatorTok{;}
\VariableTok{sent\_1\_idx} \OperatorTok{=}\NormalTok{ (}\VariableTok{transmitted} \OperatorTok{==} \FloatTok{1}\NormalTok{)}\OperatorTok{;}

\CommentTok{\% P(received = 1 | sent = 0)}
\VariableTok{p\_1\_given\_sent0} \OperatorTok{=} \VariableTok{mean}\NormalTok{(}\VariableTok{received}\NormalTok{(}\VariableTok{sent\_0\_idx}\NormalTok{))}\OperatorTok{;}
\CommentTok{\% P(received = 0 | sent = 1)  }
\VariableTok{p\_0\_given\_sent1} \OperatorTok{=} \VariableTok{mean}\NormalTok{(}\VariableTok{received}\NormalTok{(}\VariableTok{sent\_1\_idx}\NormalTok{))}\OperatorTok{;}

\VariableTok{fprintf}\NormalTok{(}\SpecialStringTok{\textquotesingle{}P(receive 1 | sent 0) = \%.4f (theory: \%.4f)\textbackslash{}n\textquotesingle{}}\OperatorTok{,} \VariableTok{p\_1\_given\_sent0}\OperatorTok{,} \VariableTok{p\_crossover}\NormalTok{)}\OperatorTok{;}
\VariableTok{fprintf}\NormalTok{(}\SpecialStringTok{\textquotesingle{}P(receive 0 | sent 1) = \%.4f (theory: \%.4f)\textbackslash{}n\textquotesingle{}}\OperatorTok{,} \VariableTok{p\_0\_given\_sent1}\OperatorTok{,} \VariableTok{p\_crossover}\NormalTok{)}\OperatorTok{;}
\end{Highlighting}
\end{Shaded}

\end{PSbox}

\PSsubsection{مثال 4: شبیه‌سازی قضیه بیز}

\tcbset{PScodemode/.style={unbreakable}}

\begin{Shaded}
\begin{Highlighting}[]
\CommentTok{\%\% Bayes\textquotesingle{} Theorem: Fault Diagnosis Simulation}
\VariableTok{N} \OperatorTok{=} \FloatTok{100000}\OperatorTok{;}

\CommentTok{\% Prior probabilities of fault types}
\VariableTok{p\_power} \OperatorTok{=} \FloatTok{0.3}\OperatorTok{;}   \VariableTok{p\_proc} \OperatorTok{=} \FloatTok{0.5}\OperatorTok{;}   \VariableTok{p\_mem} \OperatorTok{=} \FloatTok{0.2}\OperatorTok{;}

\CommentTok{\% Likelihoods of symptom (overheating) given fault type}
\VariableTok{p\_sym\_power} \OperatorTok{=} \FloatTok{0.8}\OperatorTok{;}  \VariableTok{p\_sym\_proc} \OperatorTok{=} \FloatTok{0.3}\OperatorTok{;}  \VariableTok{p\_sym\_mem} \OperatorTok{=} \FloatTok{0.1}\OperatorTok{;}

\CommentTok{\% Simulate: generate fault type for each trial}
\VariableTok{fault\_type} \OperatorTok{=} \VariableTok{zeros}\NormalTok{(}\FloatTok{1}\OperatorTok{,} \VariableTok{N}\NormalTok{)}\OperatorTok{;}
\VariableTok{r} \OperatorTok{=} \VariableTok{rand}\NormalTok{(}\FloatTok{1}\OperatorTok{,} \VariableTok{N}\NormalTok{)}\OperatorTok{;}
\VariableTok{fault\_type}\NormalTok{(}\VariableTok{r} \OperatorTok{\textless{}} \VariableTok{p\_power}\NormalTok{) }\OperatorTok{=} \FloatTok{1}\OperatorTok{;}                          \CommentTok{\% Power supply}
\VariableTok{fault\_type}\NormalTok{(}\VariableTok{r} \OperatorTok{\textgreater{}=} \VariableTok{p\_power} \OperatorTok{\&} \VariableTok{r} \OperatorTok{\textless{}} \VariableTok{p\_power} \OperatorTok{+} \VariableTok{p\_proc}\NormalTok{) }\OperatorTok{=} \FloatTok{2}\OperatorTok{;}  \CommentTok{\% Processor}
\VariableTok{fault\_type}\NormalTok{(}\VariableTok{r} \OperatorTok{\textgreater{}=} \VariableTok{p\_power} \OperatorTok{+} \VariableTok{p\_proc}\NormalTok{) }\OperatorTok{=} \FloatTok{3}\OperatorTok{;}                \CommentTok{\% Memory}

\CommentTok{\% Generate symptom based on fault type}
\VariableTok{symptom} \OperatorTok{=} \VariableTok{zeros}\NormalTok{(}\FloatTok{1}\OperatorTok{,} \VariableTok{N}\NormalTok{)}\OperatorTok{;}
\VariableTok{symptom}\NormalTok{(}\VariableTok{fault\_type} \OperatorTok{==} \FloatTok{1}\NormalTok{) }\OperatorTok{=} \VariableTok{rand}\NormalTok{(}\FloatTok{1}\OperatorTok{,} \VariableTok{sum}\NormalTok{(}\VariableTok{fault\_type}\OperatorTok{==}\FloatTok{1}\NormalTok{)) }\OperatorTok{\textless{}} \VariableTok{p\_sym\_power}\OperatorTok{;}
\VariableTok{symptom}\NormalTok{(}\VariableTok{fault\_type} \OperatorTok{==} \FloatTok{2}\NormalTok{) }\OperatorTok{=} \VariableTok{rand}\NormalTok{(}\FloatTok{1}\OperatorTok{,} \VariableTok{sum}\NormalTok{(}\VariableTok{fault\_type}\OperatorTok{==}\FloatTok{2}\NormalTok{)) }\OperatorTok{\textless{}} \VariableTok{p\_sym\_proc}\OperatorTok{;}
\VariableTok{symptom}\NormalTok{(}\VariableTok{fault\_type} \OperatorTok{==} \FloatTok{3}\NormalTok{) }\OperatorTok{=} \VariableTok{rand}\NormalTok{(}\FloatTok{1}\OperatorTok{,} \VariableTok{sum}\NormalTok{(}\VariableTok{fault\_type}\OperatorTok{==}\FloatTok{3}\NormalTok{)) }\OperatorTok{\textless{}} \VariableTok{p\_sym\_mem}\OperatorTok{;}

\CommentTok{\% Compute posterior: P(fault\_type | symptom = 1)}
\VariableTok{sym\_idx} \OperatorTok{=} \VariableTok{logical}\NormalTok{(}\VariableTok{symptom}\NormalTok{)}\OperatorTok{;}
\VariableTok{p\_power\_given\_sym} \OperatorTok{=} \VariableTok{mean}\NormalTok{(}\VariableTok{fault\_type}\NormalTok{(}\VariableTok{sym\_idx}\NormalTok{) }\OperatorTok{==} \FloatTok{1}\NormalTok{)}\OperatorTok{;}
\VariableTok{p\_proc\_given\_sym} \OperatorTok{=} \VariableTok{mean}\NormalTok{(}\VariableTok{fault\_type}\NormalTok{(}\VariableTok{sym\_idx}\NormalTok{) }\OperatorTok{==} \FloatTok{2}\NormalTok{)}\OperatorTok{;}
\VariableTok{p\_mem\_given\_sym} \OperatorTok{=} \VariableTok{mean}\NormalTok{(}\VariableTok{fault\_type}\NormalTok{(}\VariableTok{sym\_idx}\NormalTok{) }\OperatorTok{==} \FloatTok{3}\NormalTok{)}\OperatorTok{;}

\CommentTok{\% Theoretical values}
\VariableTok{p\_sym} \OperatorTok{=} \VariableTok{p\_sym\_power}\OperatorTok{*}\VariableTok{p\_power} \OperatorTok{+} \VariableTok{p\_sym\_proc}\OperatorTok{*}\VariableTok{p\_proc} \OperatorTok{+} \VariableTok{p\_sym\_mem}\OperatorTok{*}\VariableTok{p\_mem}\OperatorTok{;}
\VariableTok{theory\_power} \OperatorTok{=} \VariableTok{p\_sym\_power}\OperatorTok{*}\VariableTok{p\_power} \OperatorTok{/} \VariableTok{p\_sym}\OperatorTok{;}
\VariableTok{theory\_proc} \OperatorTok{=} \VariableTok{p\_sym\_proc}\OperatorTok{*}\VariableTok{p\_proc} \OperatorTok{/} \VariableTok{p\_sym}\OperatorTok{;}
\VariableTok{theory\_mem} \OperatorTok{=} \VariableTok{p\_sym\_mem}\OperatorTok{*}\VariableTok{p\_mem} \OperatorTok{/} \VariableTok{p\_sym}\OperatorTok{;}

\VariableTok{fprintf}\NormalTok{(}\SpecialStringTok{\textquotesingle{}P(power fault | symptom): Empirical = \%.3f, Theory = \%.3f\textbackslash{}n\textquotesingle{}}\OperatorTok{,} \VariableTok{p\_power\_given\_sym}\OperatorTok{,} \VariableTok{theory\_power}\NormalTok{)}\OperatorTok{;}
\VariableTok{fprintf}\NormalTok{(}\SpecialStringTok{\textquotesingle{}P(proc fault  | symptom): Empirical = \%.3f, Theory = \%.3f\textbackslash{}n\textquotesingle{}}\OperatorTok{,} \VariableTok{p\_proc\_given\_sym}\OperatorTok{,} \VariableTok{theory\_proc}\NormalTok{)}\OperatorTok{;}
\VariableTok{fprintf}\NormalTok{(}\SpecialStringTok{\textquotesingle{}P(mem fault   | symptom): Empirical = \%.3f, Theory = \%.3f\textbackslash{}n\textquotesingle{}}\OperatorTok{,} \VariableTok{p\_mem\_given\_sym}\OperatorTok{,} \VariableTok{theory\_mem}\NormalTok{)}\OperatorTok{;}
\end{Highlighting}
\end{Shaded}

\begin{PSbox}[unbreakable]{ex}{مثال 5: آزمون استقلال}

\tcbset{PScodemode/.style={unbreakable}}

\begin{Shaded}
\begin{Highlighting}[]
\CommentTok{\%\% Testing Independence vs. Dependence}
\VariableTok{N} \OperatorTok{=} \FloatTok{100000}\OperatorTok{;}

\CommentTok{\% Case 1: Independent events}
\VariableTok{A\_ind} \OperatorTok{=} \VariableTok{rand}\NormalTok{(}\FloatTok{1}\OperatorTok{,} \VariableTok{N}\NormalTok{) }\OperatorTok{\textless{}} \FloatTok{0.3}\OperatorTok{;}
\VariableTok{B\_ind} \OperatorTok{=} \VariableTok{rand}\NormalTok{(}\FloatTok{1}\OperatorTok{,} \VariableTok{N}\NormalTok{) }\OperatorTok{\textless{}} \FloatTok{0.4}\OperatorTok{;}  \CommentTok{\% Generated independently}

\VariableTok{pA} \OperatorTok{=} \VariableTok{mean}\NormalTok{(}\VariableTok{A\_ind}\NormalTok{)}\OperatorTok{;}
\VariableTok{pB} \OperatorTok{=} \VariableTok{mean}\NormalTok{(}\VariableTok{B\_ind}\NormalTok{)}\OperatorTok{;}
\VariableTok{pAB} \OperatorTok{=} \VariableTok{mean}\NormalTok{(}\VariableTok{A\_ind} \OperatorTok{\&} \VariableTok{B\_ind}\NormalTok{)}\OperatorTok{;}
\VariableTok{fprintf}\NormalTok{(}\SpecialStringTok{\textquotesingle{}Independent case:\textbackslash{}n\textquotesingle{}}\NormalTok{)}\OperatorTok{;}
\VariableTok{fprintf}\NormalTok{(}\SpecialStringTok{\textquotesingle{}  P(A)*P(B) = \%.4f, P(A∩B) = \%.4f\textbackslash{}n\textquotesingle{}}\OperatorTok{,} \VariableTok{pA}\OperatorTok{*}\VariableTok{pB}\OperatorTok{,} \VariableTok{pAB}\NormalTok{)}\OperatorTok{;}

\CommentTok{\% Case 2: Dependent events (B depends on A)}
\VariableTok{A\_dep} \OperatorTok{=} \VariableTok{rand}\NormalTok{(}\FloatTok{1}\OperatorTok{,} \VariableTok{N}\NormalTok{) }\OperatorTok{\textless{}} \FloatTok{0.3}\OperatorTok{;}
\VariableTok{B\_dep} \OperatorTok{=} \VariableTok{zeros}\NormalTok{(}\FloatTok{1}\OperatorTok{,} \VariableTok{N}\NormalTok{)}\OperatorTok{;}
\VariableTok{B\_dep}\NormalTok{(}\VariableTok{A\_dep} \OperatorTok{==} \FloatTok{1}\NormalTok{) }\OperatorTok{=} \VariableTok{rand}\NormalTok{(}\FloatTok{1}\OperatorTok{,} \VariableTok{sum}\NormalTok{(}\VariableTok{A\_dep}\NormalTok{)) }\OperatorTok{\textless{}} \FloatTok{0.7}\OperatorTok{;}   \CommentTok{\% P(B|A) = 0.7}
\VariableTok{B\_dep}\NormalTok{(}\VariableTok{A\_dep} \OperatorTok{==} \FloatTok{0}\NormalTok{) }\OperatorTok{=} \VariableTok{rand}\NormalTok{(}\FloatTok{1}\OperatorTok{,} \VariableTok{sum}\NormalTok{(}\OperatorTok{\textasciitilde{}}\VariableTok{A\_dep}\NormalTok{)) }\OperatorTok{\textless{}} \FloatTok{0.2}\OperatorTok{;}  \CommentTok{\% P(B|not A) = 0.2}

\VariableTok{pA} \OperatorTok{=} \VariableTok{mean}\NormalTok{(}\VariableTok{A\_dep}\NormalTok{)}\OperatorTok{;}
\VariableTok{pB} \OperatorTok{=} \VariableTok{mean}\NormalTok{(}\VariableTok{B\_dep}\NormalTok{)}\OperatorTok{;}
\VariableTok{pAB} \OperatorTok{=} \VariableTok{mean}\NormalTok{(}\VariableTok{A\_dep} \OperatorTok{\&} \VariableTok{B\_dep}\NormalTok{)}\OperatorTok{;}
\VariableTok{fprintf}\NormalTok{(}\SpecialStringTok{\textquotesingle{}Dependent case:\textbackslash{}n\textquotesingle{}}\NormalTok{)}\OperatorTok{;}
\VariableTok{fprintf}\NormalTok{(}\SpecialStringTok{\textquotesingle{}  P(A)*P(B) = \%.4f, P(A∩B) = \%.4f\textbackslash{}n\textquotesingle{}}\OperatorTok{,} \VariableTok{pA}\OperatorTok{*}\VariableTok{pB}\OperatorTok{,} \VariableTok{pAB}\NormalTok{)}\OperatorTok{;}
\VariableTok{fprintf}\NormalTok{(}\SpecialStringTok{\textquotesingle{}  These differ =\textgreater{} events are dependent\textbackslash{}n\textquotesingle{}}\NormalTok{)}\OperatorTok{;}
\end{Highlighting}
\end{Shaded}

\end{PSbox}

\PSsection{1.12}{خلاصه: تمایزهای کلیدی}

{\def\LTcaptype{none} % do not increment counter
\begin{longtable}[c]{@{}
  >{\centering\arraybackslash}p{(0.965\linewidth - 4\tabcolsep) * \real{0.2195}}
  >{\centering\arraybackslash}p{(0.965\linewidth - 4\tabcolsep) * \real{0.2683}}
  >{\centering\arraybackslash}p{(0.965\linewidth - 4\tabcolsep) * \real{0.5122}}@{}}
\toprule\noalign{}
\rowcolor{PSnavy}\PSth{مفهوم} & \PSth{تعریف} & \PSth{معنای مهندسی} \\
\midrule\noalign{}
\endhead
\bottomrule\noalign{}
\endlastfoot
\textbf{رویداد} & زیرمجموعه‌ای از فضای نمونه & یک خروجی مشخص یا مجموعه‌ای از خروجی‌ها که مورد توجه ما هستند \\
\textbf{احتمال} & \lr{$P(A)$} = معیار احتمال وقوع & فراوانی نسبی بلندمدت وقوع \\
\textbf{احتمال شرطی} & \lr{$P(A \mid B)$} = احتمال به‌روزشده با توجه به \lr{$B$} & اینکه دانستن یک رویداد چگونه احتمال رویداد دیگر را تغییر می‌دهد \\
\textbf{استقلال} & \lr{$P(A \cap B) = P(A)P(B)$} & دانستن یک رویداد هیچ اطلاعاتی درباره دیگری نمی‌دهد \\
\end{longtable}
}

\PSsection{1.13}{مسایل تمرینی}

\begin{PSbox}[unbreakable]{prob}{مسیله 1: قابلیت اطمینان قطعات}

یک سیستم دارای دو پردازنده مستقل است. پردازنده 1 با احتمال \PSnum{0.02} و پردازنده 2 با احتمال \PSnum{0.03} خراب می‌شود. سیستم فقط در صورتی از کار می‌افتد که هر دو پردازنده خراب شوند.

\begin{enumerate}
\def\labelenumi{(\PSalph{enumi})}
\tightlist
\item
  احتمال خرابی سیستم چقدر است؟
\item
  احتمال اینکه دست‌کم یکی از پردازنده‌ها خراب شود چقدر است؟
\item
  با توجه به اینکه سیستم همچنان در حال کار است، احتمال اینکه پردازنده 1 خراب شده باشد چقدر است؟
\end{enumerate}

\textbf{راه‌حل:}

\begin{enumerate}
\def\labelenumi{(\PSalph{enumi})}
\item
  \lr{$P(\text{\rl{خرابی سیستم}}) = P(\text{\rl{هر دو خراب شوند}}) = (0.02)(0.03) = 0.0006$}
\item
  \lr{$P(\text{\rl{دست‌کم یکی خراب شود}}) = 1 - P(\text{\rl{هیچ‌کدام خراب نشوند}}) = 1 - (0.98)(0.97) = 1 - 0.9506 = 0.0494$}
\item
  سیستم در حال کار = \lr{\mbox{NOT(\rl{هر دو خراب شوند})}} = دست‌کم یکی کار می‌کند.\PSbr{}\lr{$\displaystyle P(P1\ \PSmtext{failed}\ \mid \PSmtext{system works}) = \frac{P(P1\ \PSmtext{failed}\ \cap P2\ \PSmtext{works})}{P(\PSmtext{system works})}$}\PSbr{}\lr{$\displaystyle = \frac{(0.02)(0.97)}{1 - 0.0006} = \frac{0.0194}{0.9994} \approx 0.0194$}
\end{enumerate}

\end{PSbox}

\begin{PSbox}[unbreakable]{prob}{مسیله 2: آشکارسازی حسگر}

یک سیستم آشکارسازی آتش دارای مقادیر زیر است: \lr{$P(\PSmtext{fire}) = 0.001$}، \lr{$P(\PSmtext{alarm}\ \mid \PSmtext{fire}) = 0.98$}، \lr{$P(\PSmtext{alarm}\ \mid \PSmtext{no fire}) = 0.01$}.

\begin{enumerate}
\def\labelenumi{(\PSalph{enumi})}
\tightlist
\item
  \lr{$P(\PSmtext{alarm})$} چقدر است؟
\item
  اگر هشدار به صدا درآید، \lr{$P(\PSmtext{fire}\ \mid \PSmtext{alarm})$} چقدر است؟
\item
  نتیجه را تفسیر کنید.
\end{enumerate}

\textbf{راه‌حل:}

\begin{enumerate}
\def\labelenumi{(\PSalph{enumi})}
\item
  \lr{$P(\PSmtext{alarm}) = P(\PSmtext{alarm}\ \mid \PSmtext{fire})P(\PSmtext{fire}) + P(\PSmtext{alarm}\ \mid \PSmtext{no fire})P(\PSmtext{no fire})$}\PSbr{}\lr{$= (0.98)(0.001) + (0.01)(0.999) = 0.000980 + 0.009990 = 0.010970$}
\item
  \lr{$\displaystyle P(\PSmtext{fire}\ \mid \PSmtext{alarm}) = \frac{P(\PSmtext{alarm}\ \mid \PSmtext{fire})P(\PSmtext{fire})}{P(\PSmtext{alarm})}$}\PSbr{}\lr{$\displaystyle = \frac{(0.98)(0.001)}{0.01097} = 0.0894$}
\item
  فقط حدود 9\% از هشدارها مربوط به آتش‌سوزی واقعی هستند! این همان مغالطه نرخ پایه است. حتی با وجود یک آشکارساز عالی، رویدادهای نادر هشدارهای کاذب زیادی ایجاد می‌کنند.
\end{enumerate}

\end{PSbox}

\begin{PSbox}[unbreakable]{prob}{مسیله 3: خطای ارتباطی}

یک بسته 8 بیتی ارسال می‌شود. هر بیت دارای احتمال خطای مستقل \PSnum{0.05} است.

\begin{enumerate}
\def\labelenumi{(\PSalph{enumi})}
\tightlist
\item
  \lr{$P(\text{\rl{عدم وجود خطا در بسته}})$} چقدر است؟
\item
  \lr{$P(\text{\rl{دقیقاً یک خطا}})$} چقدر است؟
\item
  \lr{$P(\text{\rl{دست‌کم یک خطا}})$} چقدر است؟
\item
  کد \lr{\mbox{MATLAB}} را برای راستی‌آزمایی از طریق شبیه‌سازی بنویسید.
\end{enumerate}

\textbf{راه‌حل:}

\begin{enumerate}
\def\labelenumi{(\PSalph{enumi})}
\item
  \lr{$P(\text{\rl{عدم وجود خطا}}) = (0.95)^{8} = 0.6634$}
\item
  \lr{$P(\text{\rl{دقیقاً یک خطا}}) = C(8,\,1)(0.05)^{1}(0.95)^{7} = 8 \times 0.05 \times 0.6983 = 0.2793$}
\item
  \lr{$P(\text{\rl{دست‌کم یک خطا}}) = 1 - P(\text{\rl{عدم وجود خطا}}) = 1 - 0.6634 = 0.3366$}
\item
  راستی‌آزمایی \lr{\mbox{MATLAB}}:
\end{enumerate}

\tcbset{PScodemode/.style={unbreakable}}

\begin{Shaded}
\begin{Highlighting}[]
\VariableTok{p} \OperatorTok{=} \FloatTok{0.05}\OperatorTok{;} \VariableTok{n\_bits} \OperatorTok{=} \FloatTok{8}\OperatorTok{;} \VariableTok{N\_packets} \OperatorTok{=} \FloatTok{100000}\OperatorTok{;}
\VariableTok{errors\_per\_packet} \OperatorTok{=} \VariableTok{sum}\NormalTok{(}\VariableTok{rand}\NormalTok{(}\VariableTok{n\_bits}\OperatorTok{,} \VariableTok{N\_packets}\NormalTok{) }\OperatorTok{\textless{}} \VariableTok{p}\OperatorTok{,} \FloatTok{1}\NormalTok{)}\OperatorTok{;}
\VariableTok{fprintf}\NormalTok{(}\SpecialStringTok{\textquotesingle{}P(0 errors) = \%.4f (theory: \%.4f)\textbackslash{}n\textquotesingle{}}\OperatorTok{,} \VariableTok{mean}\NormalTok{(}\VariableTok{errors\_per\_packet}\OperatorTok{==}\FloatTok{0}\NormalTok{)}\OperatorTok{,} \FloatTok{0.95}\OperatorTok{\^{}}\FloatTok{8}\NormalTok{)}\OperatorTok{;}
\VariableTok{fprintf}\NormalTok{(}\SpecialStringTok{\textquotesingle{}P(1 error)  = \%.4f (theory: \%.4f)\textbackslash{}n\textquotesingle{}}\OperatorTok{,} \VariableTok{mean}\NormalTok{(}\VariableTok{errors\_per\_packet}\OperatorTok{==}\FloatTok{1}\NormalTok{)}\OperatorTok{,} \FloatTok{8}\OperatorTok{*}\FloatTok{0.05}\OperatorTok{*}\FloatTok{0.95}\OperatorTok{\^{}}\FloatTok{7}\NormalTok{)}\OperatorTok{;}
\VariableTok{fprintf}\NormalTok{(}\SpecialStringTok{\textquotesingle{}P(\textgreater{}=1 error)= \%.4f (theory: \%.4f)\textbackslash{}n\textquotesingle{}}\OperatorTok{,} \VariableTok{mean}\NormalTok{(}\VariableTok{errors\_per\_packet}\OperatorTok{\textgreater{}=}\FloatTok{1}\NormalTok{)}\OperatorTok{,} \FloatTok{1}\OperatorTok{{-}}\FloatTok{0.95}\OperatorTok{\^{}}\FloatTok{8}\NormalTok{)}\OperatorTok{;}
\end{Highlighting}
\end{Shaded}

\end{PSbox}

\begin{PSbox}[unbreakable]{prob}{مسیله 4: کاربرد احتمال کل}

یک سیستم بی‌سیم در سه حالت کار می‌کند:

\begin{itemize}
\tightlist
\item
  حالت 1 (کانال خوب): 60\% از زمان، \lr{$P(\PSmtext{outage}) = 0.01$}
\item
  حالت 2 (کانال متوسط): 30\% از زمان، \lr{$P(\PSmtext{outage}) = 0.10$}
\item
  حالت 3 (کانال ضعیف): 10\% از زمان، \lr{$P(\PSmtext{outage}) = 0.50$}
\end{itemize}

\begin{enumerate}
\def\labelenumi{(\PSalph{enumi})}
\tightlist
\item
  احتمال کلی قطعی چقدر است؟
\item
  با توجه به اینکه قطعی رخ داده است، احتمال اینکه سیستم در حالت 3 بوده باشد چقدر است؟
\end{enumerate}

\textbf{راه‌حل:}

\begin{enumerate}
\def\labelenumi{(\PSalph{enumi})}
\item
  \lr{$P(\PSmtext{outage}) = (0.01)(0.60) + (0.10)(0.30) + (0.50)(0.10) = 0.006 + 0.030 + 0.050 = 0.086$}
\item
  \lr{$\displaystyle P(\PSmtext{Mode 3}\ \mid \PSmtext{outage}) = \frac{(0.50)(0.10)}{0.086} = 0.581$}
\end{enumerate}

\end{PSbox}

\begin{PSbox}[unbreakable]{prob}{مسیله 5: تحلیل استقلال}

دو حسگر یک فرایند را پایش می‌کنند. حسگر \lr{$A$} دارای احتمال هشدار کاذب \PSnum{0.05} است. حسگر \lr{$B$} دارای احتمال هشدار کاذب \PSnum{0.03} است.

\begin{enumerate}
\def\labelenumi{(\PSalph{enumi})}
\tightlist
\item
  اگر حسگرها مستقل باشند، \lr{$P(\text{\rl{هر دو هشدار کاذب می‌دهند}})$} چقدر است؟
\item
  اگر حسگرها یک منبع نویز مشترک داشته باشند و \lr{$P(\text{\rl{هر دو هشدار کاذب}}) = 0.01$} باشد، آیا مستقل هستند؟
\item
  در حالت \lr{\mbox{(b)}}، \lr{$P(A\ \PSmtext{false alarm}\ \mid B\ \PSmtext{false alarm})$} چقدر است؟
\end{enumerate}

\textbf{راه‌حل:}

\begin{enumerate}
\def\labelenumi{(\PSalph{enumi})}
\item
  در صورت استقلال: \lr{$P(\text{\rl{هر دو}}) = (0.05)(0.03) = 0.0015$}
\item
  اگر \lr{$P(\text{\rl{هر دو}}) = 0.01 \ne 0.0015$}، آن‌ها مستقل نیستند. منبع نویز مشترک وابستگی ایجاد می‌کند.
\item
  \lr{$\displaystyle P(A \mid B) = \frac{P(A \cap B)}{P(B)} = \frac{0.01}{0.03} = 0.333$}
\end{enumerate}

این مقدار بسیار بیشتر از \lr{$P(A) = 0.05$} است و وابستگی شدید را تأیید می‌کند.

\end{PSbox}

\PSsection{1.14}{نکات کلیدی}

\begin{enumerate}
\def\labelenumi{\arabic{enumi}.}
\tightlist
\item
  \textbf{احتمال عدم‌قطعیت را کمّی می‌کند} \lr{—} برای طراحی مهندسی تحت عدم‌قطعیت ضروری است
\item
  \textbf{احتمال شرطی باورها را به‌روزرسانی می‌کند} \lr{—} مبنای آشکارسازی، تخمین و تشخیص
\item
  \textbf{استقلال تحلیل را ساده می‌کند} \lr{—} اما باید توجیه شود، نه اینکه صرفاً فرض شود
\item
  \textbf{قضیه بیز احتمال‌های شرطی را معکوس می‌کند} \lr{—} از «احتمال شواهد با توجه به علت» به «احتمال علت با توجه به شواهد»
\item
  \textbf{نرخ‌های پایه اهمیت دارند} \lr{—} حتی آشکارسازهای خوب نیز هنگامی که رویدادها نادر هستند، هشدار کاذب تولید می‌کنند
\item
  \textbf{شبیه‌سازی تحلیل را اعتبارسنجی می‌کند} \lr{—} مونت‌کارلوی \lr{\mbox{MATLAB}} محاسبات نظری را تأیید می‌کند
\end{enumerate}

\emph{ماژول بعدی: {ماژول 2 \lr{—} یک متغیر تصادفی و توابع آن}}
